\section{The polar obstruction without a zero-free half-plane premise}
\label{sec:polar-critical-line-obstruction}

The rank-two obstruction can be strengthened independently of
the full-matrix extension. The known critical-line zeros alone
already make the two polar quotient classes independent.
This uses a proved partial zero-location theorem, not A, M,
or a supposition that the whole divisor lies in the strip.

Fix any \(0<a<1/2\), \(s\ge0\), and the actual rational
completion \(\mathcal W=H_{\Sigma,a,s}\oplus\mathbb C^2\)
and \(\mathcal W_+=J\oplus\mathbb C^2\) of MA-MO-016.
The sections, cutoff and functions \(r_0,r_1\) are unchanged;
write \(G=\mathcal Mg\) and \(P=\mathcal M\phi'\).

\paragraph{MA-MO-023 (unconditional polar quotient independence).}
For every nonzero \(g\in C_c^\infty(\mathbb R)\),
\[
 \mathbb C^2\longrightarrow H_{\Sigma,a,s}/J,
       \qquad c\longmapsto[C_g(c_0r_0+c_1r_1)]
 \tag{MA-MO.500}
\]
is injective, without (MA-MO.320). Consequently the original
smooth positive source column of \(C_g(P_+-P_-)\) has an
outgoing quotient map of rank exactly two and does not take
all its source vectors into \(\mathcal W_+\).
Whenever that column has the bounded trace-class extension
of MA-MO-016, its outgoing rank is also exactly two and the
minimum correction rank of MA-MO.403 is two. This includes
the unconditional generic-width Gevrey regime. The same
statements hold when \(g(0)=G(1/2)=G(-1/2)=0\), as long
as \(g\ne0\). No unrestricted full-matrix extension is
asserted here.

\begin{proof}
For every real \(\gamma\), the centered Mellin evaluation
\[
              f\longmapsto\int f(x)e^{i\gamma x}\,dx
 \tag{MA-MO.501}
\]
is a continuous functional on \(H_{\Sigma,a,s}\).
Indeed its squared dual norm is
\(\int[(e^{2ax}+e^{-2ax})(1+x^2)^{s/2}]^{-1}dx<\infty\).
If \(\rho=1/2+i\gamma\) is a nontrivial zeta zero, this
functional annihilates the smooth theta ideal and therefore
its closure \(J\). This follows from the actual Mellin
theta identity, already used in MA-MO.407, with the two
source moments zero. Only a value constraint is needed below.

The exact residue identities (MA-MO.404)--(MA-MO.408)
are unconditional identities of the original smooth source.
For \(c\ne0\), their entire numerator
\[
 E_c(z)=-c_0(z-1/2)P(z+1/2)+c_1(z+1/2)P(z-1/2)
 \tag{MA-MO.502}
\]
is nonzero because \(E_c(1/2)=-c_1\) and
\(E_c(-1/2)=-c_0\). If the class in (MA-MO.500)
were zero, (MA-MO.501) and these identities would force
\[
                    G(i\gamma)E_c(i\gamma)=0
 \tag{MA-MO.503}
\]
at every critical-line zero. The denominators
\(z\pm1/2\) are nonzero at these points, so this
implication entails no endpoint division.

For completeness, specify precisely the known counting
input. Conrey (1989), Theorem 1, printed p.~4, with the
counts in equations (8), (11), and (12), printed p.~3,
implies that there are constants \(c>0\) and \(T_0\)
such that the number of simple critical-line zeros with
\(0<\gamma<T\) is at least \(cT\log T\) for all
\(T\ge T_0\). Only this weaker consequence is used.
They are distinct points \(i\gamma\), so (MA-MO.503)
forces that many distinct zeros of \(GE_c\).

But \(g\ne0\) makes \(G\) nonzero, and compact support
of \(g\) and \(\phi'\) gives the exponential-type bound
(MA-MO.411). The Jensen argument (MA-MO.412), which does
not use the half-plane premise, therefore bounds its number
of zeros in \(|z|\le T\), even with multiplicity, by
\(O(T)\). This contradicts the preceding lower bound.
Thus (MA-MO.500) is injective.

On the original positive source image, the algebraic
outgoing formula (MA-MO.256) holds before any analytic
extension is invoked. Every polar coordinate in
\(\mathbb C^2\) occurs there, by the actual two Gaussian
sections. The just-proved injectivity therefore gives rank
two already on this smooth image. If the bounded extension
exists, the same outgoing formula and density preserve that
rank. Factoring a correction through the quotient gives
the lower bound two on its rank; the orthogonal complement
realizes it exactly, as in the final paragraph of MA-MO-020.
None of these steps uses the value \(g(0)\) or either
polar Mellin value \(G(1/2),G(-1/2)\), so the stated
subclasses follow.
\end{proof}

\paragraph{Source and zero-location scope.}
Conrey's theorem is a known published input, cited rather
than reproved. Its statement, counting conventions and
locators were read directly in the original article; this
is not a new audit of its complete mollifier proof. The
unconditional lower bound concerns an established subset
of zeros. It does not say where the other zeros lie and
does not imply their simplicity. Thus MA-MO-023 removes
the containing-half-plane premise from the obstruction,
not from the extension or prime-trace evaluation in
MA-MO-018/021.

Conrey, J. B. (1989). More than two fifths of the zeros
of the Riemann zeta function are on the critical line.
\emph{Journal f\"ur die reine und angewandte Mathematik,
399}, 1--26. \url{https://doi.org/10.1515/crll.1989.399.1}.

In the inherited cote 1 framework, scans 079--080 and
111--112 through BR-TR-01/02, the obstruction already
appears in the value factors of known retained local
zero modules. Allowing their full nilpotent jet structure
cannot remove it. The correction still changes the
actual analytic source map; no positivity or motivic
weight assertion is supplied by that change.
